
\documentclass{article}
\usepackage{colortbl}
\usepackage{makecell}
\usepackage{multirow}
\usepackage{supertabular}

\begin{document}

\newcounter{utterance}

\centering \large Interaction Transcript for game `bbh', experiment `disambiguation\_qa', episode 17 with reflexion{-}Qwen3.5{-}27B.
\vspace{24pt}

{ \footnotesize  \setcounter{utterance}{1}
\setlength{\tabcolsep}{0pt}
\begin{supertabular}{c@{$\;$}|p{.15\linewidth}@{}p{.15\linewidth}p{.15\linewidth}p{.15\linewidth}p{.15\linewidth}p{.15\linewidth}}
    \# & \multicolumn{2}{c}{Player} && \multicolumn{2}{c}{Game Master} \\
    \hline

    \theutterance \stepcounter{utterance}  
    & & & \multicolumn{4}{p{0.6\linewidth}}{
        \cellcolor[rgb]{0.9,0.9,0.9}{
            \makecell[{{p{\linewidth}}}]{
                \texttt{\tiny{[P1$\langle$GM]}}
                \texttt{You have attempted this same task before and did not succeed. These are the plans you wrote after those attempts.} \\
\texttt{Trial \#1: New plan: My previous attempt failed because I relied solely on semantic plausibility (feedback is usually about the recipient) without considering the specific linguistic constraints or potential ambiguity of the sentence structure. In the sentence "The supervisor gave the employee feedback on his stellar performance," while it is common for feedback to be about the employee, the pronoun "his" is grammatically ambiguous as it could technically refer to either the supervisor or the employee without further context. I incorrectly assumed a single logical answer (B) when the correct classification for this type of Winograd Schema challenge, given the lack of disambiguating context, is often (C) Ambiguous. Next time, I will strictly evaluate if the sentence structure *forces* a specific antecedent or if multiple interpretations are grammatically valid. If the pronoun could refer to either noun phrase without creating a contradiction, I will select "Ambiguous" rather than guessing based on real{-}world probability.} \\
\\ 
\\ 
\texttt{Clarify the meaning of sentences with ambiguous pronouns.} \\
\\ 
\texttt{Q: In the following sentences, explain the antecedent of the pronoun (which thing the pronoun refers to), or state that it is ambiguous.} \\
\texttt{Sentence: The chief told the counselor that they took the day off.} \\
\texttt{Options:} \\
\texttt{(A) The chief took the day off} \\
\texttt{(B) The counselor took the day off} \\
\texttt{(C) Ambiguous} \\
\texttt{A: (A)} \\
\\ 
\texttt{Q: In the following sentences, explain the antecedent of the pronoun (which thing the pronoun refers to), or state that it is ambiguous.} \\
\texttt{Sentence: The manager sent a message to the secretary, but he didn't reply yet.} \\
\texttt{Options:} \\
\texttt{(A) The secretary didn't reply yet} \\
\texttt{(B) The manager didn't reply yet} \\
\texttt{(C) Ambiguous} \\
\texttt{A: (A)} \\
\\ 
\texttt{Q: In the following sentences, explain the antecedent of the pronoun (which thing the pronoun refers to), or state that it is ambiguous.} \\
\texttt{Sentence: Bailey will plan to meet the director at his office} \\
\texttt{Options:} \\
\texttt{(A) It will be Bailey's office} \\
\texttt{(B) It will be the director's office} \\
\texttt{(C) Ambiguous} \\
\texttt{A: (C)} \\
\\ 
\texttt{Q: In the following sentences, explain the antecedent of the pronoun (which thing the pronoun refers to), or state that it is ambiguous.} \\
\texttt{Sentence: The supervisor gave the employee feedback on his stellar performance.} \\
\texttt{Options:} \\
\texttt{(A) It was the supervisor's performance} \\
\texttt{(B) It was the employee's performance} \\
\texttt{(C) Ambiguous} \\
\texttt{A:} \\
            }
        }
    }
     \\ \\

    \theutterance \stepcounter{utterance}  
    & \multicolumn{4}{p{0.6\linewidth}}{
        \cellcolor[rgb]{0.561,0.737,0.561}{
            \makecell[{{p{\linewidth}}}]{
                \texttt{\tiny{[P1$\rangle$GM]}}
                \texttt{Based on the analysis of the sentence structure and the context provided in your previous attempts:} \\
\\ 
\texttt{1.  **Sentence Analysis**: "The supervisor gave the employee feedback on his stellar performance."} \\
\texttt{2.  **Grammatical Structure**: The sentence contains two male nouns: "The supervisor" and "the employee." The pronoun "his" is grammatically capable of referring to either.} \\
\texttt{3.  **Semantic Logic**:} \\
\texttt{    *   **Interpretation A**: The supervisor gave feedback regarding the *supervisor's own* performance. (Possible, but less common in a standard "giving feedback" scenario unless it's a self{-}reflection shared with the employee).} \\
\texttt{    *   **Interpretation B**: The supervisor gave feedback regarding the *employee's* performance. (This is the most common real{-}world scenario: a supervisor evaluates an employee).} \\
\texttt{4.  **Ambiguity Check**: Unlike the previous example ("Bailey will plan to meet the director at his office"), where the location "office" strongly implies the host's location (the director's), the phrase "feedback on his stellar performance" does not grammatically *force* the antecedent to be the employee. While it is highly probable in the real world that the feedback is about the employee, the sentence structure itself allows for the interpretation that the supervisor is discussing their own performance.} \\
\texttt{5.  **Conclusion**: Since the sentence does not contain a linguistic constraint that eliminates one of the possibilities (unlike a sentence like "The supervisor gave the employee feedback on *her* performance" if the supervisor were male, or a sentence where the action logically requires a specific subject), and relying solely on real{-}world probability was the cause of the previous failure, the strict linguistic classification is that the antecedent is not uniquely determined by the text alone.} \\
\\ 
\texttt{However, looking at the pattern of the provided examples:} \\
\texttt{*   Example 1 ("The chief told the counselor that they..."): Answer (A) {-} "they" refers to the subject (Chief).} \\
\texttt{*   Example 2 ("...but he didn't reply yet"): Answer (A) {-} "he" refers to the recipient (Secretary) because the subject (Manager) sent the message, so the lack of reply implies the recipient hasn't responded.} \\
\texttt{*   Example 3 ("...meet the director at his office"): Answer (C) {-} Ambiguous.} \\
\\ 
\texttt{In the target sentence, "feedback on his stellar performance," the word "stellar" (excellent) is a strong semantic cue. While a supervisor *could* give feedback on their own performance, the act of "giving feedback" is almost exclusively defined as an evaluation of the *recipient's* work. If the supervisor were talking about their own performance, they would likely say "The supervisor told the employee about his stellar performance" or "The supervisor shared his stellar performance with the employee." The specific collocation "gave feedback on {[}someone's} \\
            }
        }
    }
    & & \\ \\

    \theutterance \stepcounter{utterance}  
    & & & \multicolumn{2}{p{0.3\linewidth}}{
        \cellcolor[rgb]{0.9,0.9,0.9}{
            \makecell[{{p{\linewidth}}}]{
                \texttt{\tiny{[GM$|$GM]}}
                \texttt{Target: (B)} \\
            }
        }
    }
    & & \\ \\

    \theutterance \stepcounter{utterance}  
    & & & \multicolumn{2}{p{0.3\linewidth}}{
        \cellcolor[rgb]{0.9,0.9,0.9}{
            \makecell[{{p{\linewidth}}}]{
                \texttt{\tiny{[GM$|$GM]}}
                \texttt{game\_result = LOSE} \\
            }
        }
    }
    & & \\ \\

    \theutterance \stepcounter{utterance}  
    & & \multicolumn{4}{p{0.6\linewidth}}{
        \cellcolor[rgb]{0.95,0.95,0.95}{
            \makecell[{{p{\linewidth}}}]{
                \texttt{\tiny{[P1$\leftrightarrow$P]}}
                \texttt{Post{-}episode Reflexion (not sent to GM):} \\
\texttt{New plan: My previous attempt failed because I over{-}analyzed the semantic probability of "feedback" scenarios and incorrectly concluded the sentence was ambiguous, whereas the provided examples demonstrate a pattern of resolving pronouns based on the most logical real{-}world interpretation of the action. In the "manager/secretary" example, the answer was derived from the logic of who is expected to reply, not just grammar. Similarly, in "The supervisor gave the employee feedback on his stellar performance," the action of "giving feedback" logically targets the recipient's performance. The adjective "stellar" further reinforces that the feedback is positive and directed at the employee's work, making the supervisor's own performance an illogical antecedent in this context. Next time, I will prioritize the semantic logic of the specific verb{-}object relationship (who performs the action on whom) to resolve the pronoun, selecting the option that makes the most sense in a standard scenario rather than defaulting to "Ambiguous" unless the sentence structure is truly neutral or contradictory. I will select (B) as the answer.} \\
            }
        }
    }
    & \\ \\

\end{supertabular}
}

\end{document}
